% ACHTUNG: Für das Erstellen des Literaturverzeichnisses wird das modernere Paket biblatex
%			 in Kombination mit biber verwendet -- nicht mehr das ältere BibTex!
% 			 Bitte stellen Sie ggf. Ihre TeX-Umgebung
% 			 entsprechend ein (z.B. TeXStudio: Einstellungen --> Erzeugen --> Standard Bibliographieprogramm: biber)
%

\documentclass[
	12pt,
	BCOR=5mm,
	DIV=12,
	headinclude=on,
	footinclude=off,
	parskip=half,
	bibliography=totoc,
	listof=entryprefix,
	toc=listof,
	pointlessnumbers,
	plainfootsepline]{scrreprt}

%	Konfigurationsdatei einziehen
% !TEX root =  master.tex
% 		HYPERREF
%
\usepackage[
	hidelinks=true % keine roten Markierungen bei Links
]{hyperref}

% Zwei eigene Befehle zum Setzen von Autor und Titel. Ausserdem werden die PDF-Informationen richtig gesetzt.
\newcommand{\TitelDerArbeit}[1]{\def\DerTitelDerArbeit{#1}\hypersetup{pdftitle={#1}}}
\newcommand{\AutorDerArbeit}[1]{\def\DerAutorDerArbeit{#1}\hypersetup{pdfauthor={#1}}}
\newcommand{\Firma}[1]{\def\DerNameDerFirma{#1}}
\newcommand{\Kurs}[1]{\def\DieKursbezeichnung{#1}}


% Correct superscripts 
\usepackage{fnpct}


%		FONT AND INPUT ENCODING
%
\usepackage[T1]{fontenc}
\usepackage[utf8]{inputenc}

%		CALCULATIONS
%
\usepackage{calc} % Used for extra space below footsepline

%		LANGUAGE SETTINGS
%
\usepackage[ngerman]{babel} 	% German language
\usepackage[german=quotes]{csquotes} 	% correct quotes using \enquote{}

%\usepackage[english]{babel}   % For english language
%\usepackage{csquotes} 	% Richtiges Setzen der Anführungszeichen mit \enquote{}


%		BIBLIOGRAPHY SETTINGS
%

% Uncomment the next three lines for author-year-style with footnotes (Chicago)
\usepackage[backend=biber, autocite=footnote, style=authoryear, dashed=false]{biblatex} 	%Use Author-Year-Cites with footnotes
\AdaptNoteOpt\footcite\multfootcite   %will add  separators if footcite is called multiple consecutive times 
\AdaptNoteOpt\autocite\multautocite % will add  separators if autocite is called multiple consecutive times

% Uncomment the next line for IEEE-style 
% \usepackage[backend=biber, autocite=inline, style=ieee]{biblatex} 	% Use IEEE-Style (e.g. [1])

% Uncomment the next line for alphabetic style 
% \usepackage[backend=biber, autocite=inline, style=alphabetic]{biblatex} 	% Use alphabetic style (e.g. [TGK12])

% Uncomment the next two lines vor Harvard-Style 
%\usepackage[backend=biber, style=apa]{biblatex} 	
%\DeclareLanguageMapping{german}{german-apa}


\DefineBibliographyStrings{ngerman}{  %Change u.a. to et al. (german only!)
	andothers = {{et\,al\adddot}},
}

%%% Uncomment the following lines to support hard URL breaks in bibliography 
%\apptocmd{\UrlBreaks}{\do\f\do\m}{}{}
%\setcounter{biburllcpenalty}{9000}% Kleinbuchstaben
%\setcounter{biburlucpenalty}{9000}% Großbuchstaben


\setlength{\bibparsep}{\parskip}		%add some space between biblatex entries in the bibliography
\addbibresource{.bib}	%Add file bibliography.bib as biblatex resource

%		CORRECT PAGE MARGINS
%			ADDED BY JANIK STRACKE
\usepackage[a4paper,
left=35mm,% Left Margin: 3.5cm
right=25mm,% Right Margin: 2.5cm
top=15mm,% Top Margin to head: 1.5cm
bottom=15mm,% Bottom Margin to foot: 1.5cm
includehead,
includefoot]{geometry}

%		SET LINE SPACING TO ONE HALF
%			ADDED BY JANIK STRACKE
\usepackage[onehalfspacing]{setspace}


%		FOOTNOTES 
%
% Count footnotes over chapters
\usepackage{chngcntr}
\counterwithout{footnote}{chapter}

%	ACRONYMS
%%%
%%% WICHTIG: Installieren Sie das neueste Acronyms-Paket!!!
%%%
\makeatletter
\usepackage{acronym}
% \@ifpackagelater{acronym}{2015/03/20}
%{%
% \renewcommand*{\aclabelfont}[1]{\textbf{\textsf{\acsfont{#1}}}}
%}%
%{%
%}%
\makeatother

%		LISTINGS
\usepackage{listings}	%Format Listings properly
\renewcommand{\lstlistingname}{Quelltext} 
\renewcommand{\lstlistlistingname}{Quelltextverzeichnis}
\lstset{numbers=left,
	numberstyle=\tiny,
	captionpos=b,
	basicstyle=\ttfamily\small}


%		EXTRA PACKAGES

\usepackage{lipsum}    %Blindtext
\usepackage{graphicx} % use various graphics formats
\usepackage[german]{varioref} 	% nicer references \vref
\usepackage{caption}	%better Captions
\usepackage{booktabs} %nicer Tabs
\usepackage{array}
%\newcolumntype{P}[1]{>{\raggedright\arraybackslash}p{#1}}


%		ALGORITHMS
\usepackage{algorithm}
\usepackage{algpseudocode}
\renewcommand{\listalgorithmname}{Algorithmenverzeichnis}
\floatname{algorithm}{Algorithmus}


%		FONT SELECTION: Entweder Latin Modern oder Times / Helvetica
\usepackage{lmodern} %Latin modern font
%\usepackage{mathptmx}  %Helvetica / Times New Roman fonts (2 lines)
%\usepackage[scaled=.92]{helvet} %Helvetica / Times New Roman fonts (2 lines)


%		PAGE HEADER / FOOTER
%	    Warning: There are some redefinitions throughout the master.tex-file!  DON'T CHANGE THESE REDEFINITIONS!
\RequirePackage[automark,headsepline,footsepline]{scrpage2}
\pagestyle{scrheadings}
\renewcommand*{\pnumfont}{\upshape\sffamily}
\renewcommand*{\headfont}{\upshape\sffamily}
\renewcommand*{\footfont}{\upshape\sffamily}
\renewcommand{\chaptermarkformat}{}
\RedeclareSectionCommand[beforeskip=0pt]{chapter}
\clearscrheadfoot

\ifoot[\rule{0pt}{\ht\strutbox+\dp\strutbox}SAP SE / DHBW Mannheim]{\rule{0pt}{\ht\strutbox+\dp\strutbox}SAP SE / DHBW Mannheim}
\ofoot[\rule{0pt}{\ht\strutbox+\dp\strutbox}\pagemark]{\rule{0pt}{\ht\strutbox+\dp\strutbox}\pagemark}

\ohead{\headmark}



\begin{document}

%% BITTE GEBEN SIE HIER DEN TITEL UND DIE AUTORIN / DEN AUTOR DER ARBEIT AN!
%% DIESE INFORMATIONEN _MÜSSEN_ GESETZT SEIN, UM TITELBLATT, ABSTRACT UND
%% EIGENSTÄNDIGKEITSERKLÄRUNG AUTOMATISCH ANZUPASSEN!
\TitelDerArbeit{Möglichkeiten zum personalisierten Lernen im Studium}
\AutorDerArbeit{Andrea Engel}
\Firma{SAP SE}
\Kurs{WWI 17 SEB}

\begin{titlepage}
\begin{minipage}{\textwidth}
		\vspace{-2cm}
		\noindent \includegraphics[scale=0.14]{img/sap.png} \hfill   \includegraphics{img/logo.jpg}
\end{minipage}
\vspace{1em}
\sffamily
\begin{center}
	\textsf{\large{}Duale Hochschule Baden-W\"urttemberg\\[1.5mm] Mannheim}\\[2em]
	\textsf{\textbf{\Large{}Seminararbeit}}\\[3mm]
	\textsf{\textbf{Xpire - Dokumentation zum Projekthergang}} \\[1.5cm]
	\textsf{\textbf{\Large{}Studiengang Wirtschaftsinformatik}\\[3mm] \textsf{Studienrichtung Software Engineering}}
	
	\vspace{3em}
%	\textsf{\Large{Sperrvermerk}}
\vfill

\begin{minipage}{\textwidth}

\begin{tabbing}
	Wissenschaftlicher Betreuer: \hspace{0.85cm}\=\kill
	Verfasser/in: \> Felix Waage, Jonas Dutzi, Milena Zahn, \\ \> Yvonne Werner, Verena Richarz, Andrea Engel \\[1.5mm]
	Firma: \> SAP SE \\[1.5mm]
	Kurs: \> WWI 17 SEB\\[1.5mm]
	Studiengangsleiter: \> Prof. Dr. Sebastian Ritterbusch  \\[1.5mm]
	Modul: \> Entwicklung mobiler Applikationen \\[1.5mm]
	Dozent/-in: \> Michael Spengler \\[1.5mm]
	Bearbeitungszeitraum: \> 11.05.2020 - 31.07.2020
	
\end{tabbing}
\end{minipage}

\end{center}

\end{titlepage}

\pagenumbering{roman} % Römische Seitennummerierung
\normalfont

%--------------------------------
% Verzeichnisse - nicht benötige Verzeichnisse bitte auskommentieren / löschen.
%--------------------------------

%   Sperrvermerk
%\chapter*{}
Die nachfolgende Arbeit enthält vertrauliche Daten der:

\vspace{0.5cm}

	SAP SE
	
	Dietmar-Hopp-Allee 16
	
	69190 Walldorf
	
\vspace{0.5cm}
Der Inhalt dieser Arbeit darf weder als Ganzes noch in Auszügen Personen außerhalb des Prüfungsprozesses und des Evaluationsverfahrens zugänglich gemacht werden, sofern keine anders lautende Genehmigung der Ausbildungsstätte vorliegt. 
\cleardoublepage


% Ehrenwörtliche Erklärung ewerkl.tex einziehen
% \input{ewerkl.tex}

%	Kurzfassung
% \input{abstract}

%	Inhaltsverzeichnis
\tableofcontents

%	Abbildungsverzeichnis
%\listoffigures

%	Tabellenverzeichnis
% \listoftables

%	Listingsverzeichnis
% \lstlistoflistings

% 	Algorithmenverzeichnis
%\listofalgorithms

% 	Abkürzungsverzeichnis (siehe Datei acronyms.tex!)
% !TEX root =  master.tex

\clearpage
\chapter*{Abkürzungsverzeichnis}	
\addcontentsline{toc}{chapter}{Abkürzungsverzeichnis}


\begin{acronym}
	\acro{DHBW}{Duale Hochschule Baden-Württemberg}

\end{acronym}


%--------------------------------
% Start des Textteils der Arbeit
%--------------------------------
\clearpage
\ihead{\chaptername~\thechapter} % Neue Header-Definition
\pagenumbering{arabic}  % Arabische Seitenzahlen

%	Anleitungs-Datei anleitung.tex einziehen. Auf diese Weise sollten Sie versuchen, für jedes einzelne
% Kapitel eine eigene Datei anzulegen und mittels input-Kommando einzuziehen.

% !TEX root =  master.tex
\chapter{Einleitung}

\section{Grundidee von Xpire}

\section{Was ist eine PWA?}
% !TEX root =  master.tex
\chapter{Analyse}

% !TEX root =  master.tex
\section{Value Proposition Canvas}

% !TEX root =  master.tex
\section{Benutzerprofilanalyse}
Im Rahmen der Benutzerprofilanalyse werden die Eigenschaften der Endanwender untersucht, um die Gebrauchstauglichkeit der Anwendung sicherzustellen. Eine wesentliche
Rolle spielen dabei die persönlichen Merkmale, Fähigkeiten und Kenntnisse der Benutzer
sowie deren Einordnung in ein bestimmtes Nutzungsverhalten.\\
Da die progressive Web App von möglichst vielen Nutzern verwendet werden soll, gibt es kaum Einschränkungen, was die Zielgruppe betrifft. Die Gemeinsamkeit der Nutzer liegt im Besitz eines Kühlschrankes, dem selbstständigem Tätigen von Lebensmittel-Einkäufen sowie im Besitz eines Smartphones oder eines Computers. Daraus lässt sich ableiten, dass die Zielgruppe überwiegend volljährig ist und ein Smartphone oder einen Computer besitzt. Laut Statista besitzen aktuell 66,5 Mio Menschen in Deutschland ein Smartphone.\autocite[vgl.][]{Statista.2020} Ca. 13 Mio Menschen in Deutschland sind Minderjährig, wobei nicht alle Minderjährigen ein Smartphone besitzen. Zieht man trotzdem die Anzahl aller Minderjährigen von der Anzahl der Menschen, die in Deutschland ein Smartphone besitzen, ab, kommt man immer noch auf eine Anzahl von mindestens 49 Mio potentiellen Nutzer. Das sind  mehr als 50\% aller deutschen Staatsbürger.\\
All diese Nutzer können einen völlig unterschiedlichen Erfahrungsschatz bezüglich der Nutzung von neuen Technologien aufweisen. Von der Rentnerin, die nur selten einen Computer verwendet bis hin zum zwanzigjährigen Studierenden, der selbstständig Web Apps programmiert, können die unterschiedlichsten Ausgangsvoraussetzungen gegeben sein. Es is davon auszugehen, dass ein Großteil der Anwendet eine durchschnittliche Technikaffinität aufweist, es aber auch Anwender gibt, die keinerlei oder kaum Erfahrungen im Umgang mit Web Apps aufweisen. Demzufolge sollte die Oberfläche gleichzeitig so konzipiert sein, dass sie auch für Laien intuitiv und verständlich ist. Ein strukturierter und klarer Aufbau könnten diesem Aspekt Rechnung tragen. Schließlich sei noch zu erwähnen, dass die Web App vorerst nur für Deutschland entwickelt wird. Zu einem späteren Zeitpunkt ist die Internationalisierung der Applikation durchaus denkbar, dafür bedarf es dann einer Möglichkeit zur Auswahl unterschiedlicher Sprachen.




% !TEX root =  master.tex
\section{Anforderungsanalyse}
% !TEX root =  master.tex
\chapter{Konzeption und Entwurf}

\section{Eingesetzte Technologien}

\section{Konzeptuelles Modell}

\section{Entwicklung von Prototypen}

% !TEX root =  master.tex
\chapter{Implementierung}


% !TEX root =  master.tex
\chapter{Zusammenfassung und Ausblick}
% Was passiert, wenn man ein Produkt löscht? Ist es dann in der Datenbank geschlöscht oder wird nur ein Attribut gesetzt? Welche Auswirkungen hätten diese Möglichkeiten?


%	Literaturverzeichnis
\ihead{} % Neue Header-Definition
%\printbibliography
\cleardoublepage

% Der Anhang beginnt hier - jedes Kapitel wird alphabetisch aufgezählt. (Anhang A, B usw.)
\appendix
\ihead{\appendixname~\thechapter} % Neue Header-Definition

% appendix.tex einziehen





\end{document}
